% Brief (grammar: small multiples, one panel for each gate and record, with
%   a mark for every output a worker can cause; atlas VIII/fig:sealed-views).
% Reader question: for one record, how many different outputs can a worker
% cause under each gate?
% Claim: 9 under the gate that uses the submitted value, 5 under the gate
% that reads the committed record and releases on request, and 1 under the
% kernel's gate. That one output is the record's parity at the slot's
% deadline, so the bottom row depends only on the record's parity.
% Evidence: every output a worker can cause for each record from 0 to 3,
% with a slot of one bit and a deadline of 2, and a clock that runs to
% time 3 [internal, harbor worlds --each for each gate; sessions
% worlds-each-submitted.txt, worlds-each-on-request.txt and
% worlds-each.txt, plotted as printed, and
% tests/harbor-research/test_sealed_gate_numbers.py checks the counts in
% the picture and the caption against them].
% Must distinguish: the output the contract allows from an output the
% worker chose; a value from the time of its release; one record from the
% next.
% Grammar chosen: time on the horizontal axis of each panel and the value
% on the vertical, with a box at the left for a run with no release, so
% that a gate's row fills up as the worker's choices grow.
% Grammar rejected: a bar of the three counts, which gives 9, 5 and 1 and
% hides that the 5 differ only in time; the table itself, which is 60 rows.
% Five-second test: which row has one mark in each panel?
\begin{figure}[H]
\centering
\pdfigurehue{pdcobalt}%
\begin{tikzpicture}[pd figure,x=.43cm,y=.43cm,
    chosen/.style={fill=pderror},
    allowed/.style={fill=pdink},
    empty/.style={fill=pdinkmuted!55}]
  % Counts taken from the tables as they are plotted: the outputs for record 0.
  \ifdefined\pdviewsfirst\else
    \newcount\pdviewsfirst \newcount\pdviewssecond \newcount\pdviewskernel
  \fi
  \global\pdviewsfirst=0 \global\pdviewssecond=0 \global\pdviewskernel=0
  % A panel for each record, and a row for each gate.
  \newcommand{\pdviewx}[1]{(#1)*6.35}
  \def\pdviewrowfirst{0}\def\pdviewrowsecond{-4.1}\def\pdviewrowkernel{-8.2}
  % Every cell of every panel, as a faint dot.
  \foreach \row in {\pdviewrowfirst,\pdviewrowsecond,\pdviewrowkernel}{
    \foreach \rec in {0,...,3}{
      \fill[empty] ({\pdviewx{\rec}},\row+.5) circle[radius=.55pt];
      \foreach \t in {0,...,3}{
        \foreach \v in {0,1}{
          \fill[empty] ({\pdviewx{\rec}+1.6+\t},\row+\v) circle[radius=.55pt];
        }
      }
    }
  }
  % Each table is a recorded run, read as it was printed: the first line
  % is the command, and every remark starts with '#'.
  \pgfplotstableread[skip first n=1]{excerpts/worlds-each-submitted.txt}\pdviewtablefirst
  \pgfplotstableread[skip first n=1]{excerpts/worlds-each-on-request.txt}\pdviewtablesecond
  \pgfplotstableread[skip first n=1]{excerpts/worlds-each.txt}\pdviewtablekernel
  % One gate's outputs: #1 is its table, #2 its row and #3 its count.
  \newcommand{\pdviewsof}[3]{%
    \pgfplotstablegetrowsof{#1}
    \pgfmathtruncatemacro{\pdlastview}{\pgfplotsretval-1}
    \foreach \i in {0,...,\pdlastview}{
      \pgfplotstablegetelem{\i}{record}\of#1 \let\rec\pgfplotsretval
      \pgfplotstablegetelem{\i}{released}\of#1 \let\lf\pgfplotsretval
      \pgfplotstablegetelem{\i}{at}\of#1 \let\at\pgfplotsretval
      \pgfplotstablegetelem{\i}{value}\of#1 \let\vl\pgfplotsretval
      \ifnum\rec=0 \global\advance#3 1 \fi
      \ifnum\lf=0
        % No release: the box at the left of the panel.
        \fill[chosen] ({\pdviewx{\rec}-.36},#2+.14) rectangle ({\pdviewx{\rec}+.36},#2+.86);
      \else
        % The output the contract allows is the record's parity at the slot's deadline.
        \pgfmathtruncatemacro{\parity}{mod(\rec,2)}
        \def\pdviewstyle{chosen}
        \ifnum\at=2 \ifnum\vl=\parity \def\pdviewstyle{allowed}\fi\fi
        \fill[\pdviewstyle] ({\pdviewx{\rec}+1.6+\at-.36},#2+\vl-.36)
          rectangle ({\pdviewx{\rec}+1.6+\at+.36},#2+\vl+.36);
        % A frame as well as a colour, so the two kinds also differ in greyscale.
        \ifnum\at=2 \ifnum\vl=\parity
          \draw[draw=pdink,line width=.5pt] ({\pdviewx{\rec}+1.6+\at-.5},#2+\vl-.5)
            rectangle ({\pdviewx{\rec}+1.6+\at+.5},#2+\vl+.5);
        \fi\fi
      \fi
    }}
  \pdviewsof{\pdviewtablefirst}{\pdviewrowfirst}{\pdviewsfirst}
  \pdviewsof{\pdviewtablesecond}{\pdviewrowsecond}{\pdviewssecond}
  \pdviewsof{\pdviewtablekernel}{\pdviewrowkernel}{\pdviewskernel}
  % Labels for the rows and the panels.
  \foreach \rec in {0,...,3}{
    \node[pd label,anchor=south] at ({\pdviewx{\rec}+2.3},3.05) {record \rec};
    \foreach \t in {0,...,3}{
      \node[pd label,text=pdinkmuted,anchor=north] at ({\pdviewx{\rec}+1.6+\t},\pdviewrowkernel-.75) {\t};
    }
    \node[pd label,text=pdinkmuted,anchor=north] at ({\pdviewx{\rec}},\pdviewrowkernel-.75) {--};
  }
  \node[pd label,anchor=north] at ({\pdviewx{1.5}+2.3},\pdviewrowkernel-2.2)
    {horizontal: the time of the release, or a dash for no release\\vertical: the value released};
  \foreach \row/\name/\count in {\pdviewrowfirst/{the first gate}/\the\pdviewsfirst,
      \pdviewrowsecond/{the second gate}/\the\pdviewssecond,
      \pdviewrowkernel/{the kernel's gate}/\the\pdviewskernel}{
    \node[pd title,anchor=south west] at (-.55,\row+1.55) {\name: \count};
    \foreach \v in {0,1}{
      \node[pd label,text=pdinkmuted,anchor=east] at (-.75,\row+\v) {\v};
    }
  }
\end{tikzpicture}
\caption{\textbf{A worker that could cause 9 outputs for one record can
cause only 1.} A framed black square is the record's parity at the slot's
deadline, and a red one is an output the worker chose.
[internal, \texttt{harbor worlds -{}-each}]}
\label{fig:sealed-views}
\end{figure}
